\chapter*{Introduction Générale}

Aujourd'hui, les entreprises digitalisent leurs processus. Les entreprises digitalisent leurs processus pour gagner du temps, réduire les erreurs et améliorer l'efficacité.

Le processus Inovice IA gère les factures fournisseurs depuis leur réception jusqu’à leur validation. Le processus Inovice IA constitue une étape clé de la gestion financière. Dans de nombreuses entreprises, le processus Inovice IA reste encore manuel ou peu automatisé. Le processus Inovice IA provoque des retards, génère des erreurs et réduit la visibilité.

Dans ce contexte, le projet se déroule au sein de l’entreprise \textbf{Virtuel Dev}. Le projet consiste à créer une plateforme qui utilise l’intelligence artificielle pour automatiser le traitement des factures fournisseurs.

La solution traite plusieurs types de factures (PDF, images, XML). La solution utilise l’OCR pour extraire les informations et l’intelligence artificielle pour analyser les données, classer les charges et détecter les anomalies.

L’application utilise une architecture simple : un frontend en \textbf{React}, un backend en \textbf{FastAPI (framework Python)} et une base de données \textbf{PostgreSQL}.

Cette fonction permet :
\begin{itemize}
\item d’importer les factures des fournisseurs ;
\item extraire automatiquement les informations clés ;
\item Classer les charges de façon intelligente ;
\item détecter les anomalies (TVA, montants, doublons) ;
\item Aider à la prise de décision pour la validation.
Voici la réponse de l'utilisateur :

\end{itemize}

Le projet vise à offrir une solution simple, fiable et intelligente. La solution facilite le traitement des factures, améliore le travail des utilisateurs et garantit une meilleure traçabilité.

Nous avons structuré le rapport ainsi :

\begin{itemize}

\item Le premier chapitre présente le cadre général du projet et l’entreprise d’accueil.

\item Le deuxième chapitre présente l’analyse et la conception du système. Le deuxième chapitre expose les besoins fonctionnels et non fonctionnels, l’architecture du système et la planification des sprints.

\item Le troisième chapitre présente le Sprint 1 dédié à l’authentification, à la gestion des utilisateurs et à l’extraction des données.

\item Le quatrième chapitre présente le Sprint 2 qui se consacre au traitement et à la consultation des données extraites.

\item Le cinquième chapitre présente le Sprint 3 qui porte sur la validation financière et la traçabilité des actions.

\end{itemize}